\subsection{Dynamic-load network evaluation with OMNeT++}

To assess routing behavior under stochastic network dynamics, we constructed an OMNeT++-based network-layer simulator using Poisson request arrivals, stochastic key-pool updates, and fidelity decay abstractions. This simulation layer is used to compare routing policies under repeated traffic load, and it complements the hardware timing package rather than replacing it.

\paragraph{Role of the OMNeT++ layer.}
The OMNeT++ study answers the network-behavior question: how blocking probability, path fidelity, route efficiency, and link-usage behavior change as traffic load increases across different topologies. It should therefore be read as the dynamic-load comparison layer of QFlow, not as the primary source of hardware timing evidence.

\paragraph{Final use in the paper.}
For the main paper, the strongest OMNeT++ figures are:
(i) blocking probability versus load on Mesh-16,
(ii) blocking probability versus load on Irregular12,
(iii) fidelity CDF on Mesh-16,
(iv) fidelity CDF on Irregular12,
(v) a per-link usage heatmap comparing Key-Aware routing and the GA-style proxy on Mesh-16, and
(vi) the topology-size scalability figure.
The most compact table is the representative-rate comparison table derived from the refined OMNeT++ summary.

\paragraph{Claim discipline.}
The dynamic-load GA comparator in the current simulator should be described as a \emph{GA-style proxy} rather than the exact final software PMO-GA selector. The OMNeT++ results support claims about dynamic-load tradeoffs, route-quality behavior, and comparison against weaker baselines such as random routing, but they should not be overclaimed as universal dominance or exact software-equivalence proof.
